\documentclass[11pt,letterpaper]{article}
\usepackage[utf8]{inputenc}
\usepackage[T1]{fontenc}
\usepackage[spanish]{babel}
\usepackage[margin=1.8cm]{geometry}
\usepackage[skip=0.4em]{parskip}
\usepackage{xcolor}
\usepackage{enumitem}
\usepackage{tcolorbox}

\definecolor{primary}{HTML}{183653}
\definecolor{accent}{HTML}{2B5B84}
\definecolor{boxbg}{HTML}{F8FAFC}
\definecolor{boxborder}{HTML}{CBD5E1}

\tcbuselibrary{skins}
\tcbset{
    enhanced,
    colback=boxbg,
    colframe=boxborder,
    boxrule=0.8pt,
    arc=3pt,
    left=8pt, right=8pt, top=5pt, bottom=5pt
}

\setlist{topsep=1pt,itemsep=1pt,parsep=0pt,leftmargin=1.1em}
\pagestyle{empty}

\begin{document}

\begin{center}
    {\Large\bfseries\color{primary} Elección de Roles --- Hito 2 (Sistemas de Información)}\\[0.4em]
    {\normalsize Cada integrante escoge un rol para trabajar durante este hito.}
\end{center}

\vspace{0.1cm}
\hrule
\vspace{0.4cm}

\begin{tcolorbox}
    \textbf{\color{primary}Rol 1: Coordinador del Proyecto y Enlace con la Empresa}
    \begin{itemize}
        \item Supervisa que el equipo cumpla con las tareas y fechas de entrega.
        \item Se comunica con la empresa (Loard) para conseguir información y agendar las pruebas con usuarios.
        \item Investiga los archivos reales de la empresa (Dropbox y formatos de obra/SST) para sacar los datos reales que se usarán en el sistema.
    \end{itemize}
\end{tcolorbox}

\begin{tcolorbox}
    \textbf{\color{primary}Rol 2: Diseñador del Prototipo (UI/UX) y Pruebas con IA}
    \begin{itemize}
        \item Diseña las pantallas y el flujo interactivo navegable en Figma (o v0) para el MVP (corte de obra, alertas de seguridad SST, dossier de calidad).
        \item Documenta el uso de IA en el diseño: qué prompt se usó, qué botó la IA y qué se corrigió a mano para ajustarlo al negocio.
    \end{itemize}
\end{tcolorbox}

\begin{tcolorbox}
    \textbf{\color{primary}Rol 3: Componentes de Interfaz y Diseño Responsivo}
    \begin{itemize}
        \item Hace la lista de componentes que se repiten (botones, tablas, formularios, tarjetas) y cómo se estructurarían en código (React).
        \item Demuestra cómo se adapta la pantalla principal en versión Computador (oficina) y Celular (campo).
    \end{itemize}
\end{tcolorbox}

\begin{tcolorbox}
    \textbf{\color{primary}Rol 4: Evaluador de Usabilidad y Pruebas con Usuarios}
    \begin{itemize}
        \item Revisa reglas básicas de usabilidad (heurísticas de Nielsen) y accesibilidad (contraste de colores y navegación por teclado).
        \item Con el Rol 1, realiza las pruebas con al menos 2 personas reales y llena la tabla: \emph{Problema encontrado vs. Cambio hecho en el diseño}.
    \end{itemize}
\end{tcolorbox}

\begin{tcolorbox}
    \textbf{\color{primary}Rol 5: Base de Datos y Conexiones (APIs)}
    \begin{itemize}
        \item Diseña un diagrama simple de base de datos con solo las tablas necesarias para las pantallas (trabajadores, proyectos, insumos, etc.).
        \item Hace la lista de conexiones (APIs) entre la interfaz y el servidor (ej: consultar proyectos, guardar reporte).
    \end{itemize}
\end{tcolorbox}

\begin{tcolorbox}
    \textbf{\color{primary}Rol 6: Calidad de la Entrega y Presentación (Pitch de 10 min)}
    \begin{itemize}
        \item Revisa la rúbrica del profesor para asegurarse de que no falte nada y compila el documento/repositorio final.
        \item Prepara las diapositivas y el guion para la exposición de 10 minutos (enfocada en los problemas que resuelve y no en el login).
    \end{itemize}
\end{tcolorbox}

\end{document}
